\documentclass[11pt]{amsart}
\usepackage{amsmath,amssymb,amsthm,mathtools}
\usepackage[margin=1in]{geometry}
\usepackage{hyperref}

\theoremstyle{plain}
\newtheorem{theorem}{Theorem}[section]
\newtheorem{proposition}[theorem]{Proposition}
\newtheorem{lemma}[theorem]{Lemma}
\newtheorem{corollary}[theorem]{Corollary}
\newtheorem{conjecture}[theorem]{Conjecture}
\theoremstyle{definition}
\newtheorem{definition}[theorem]{Definition}
\newtheorem{question}[theorem]{Question}
\theoremstyle{remark}
\newtheorem{remark}[theorem]{Remark}

\DeclareMathOperator{\GL}{GL}
\DeclareMathOperator{\SL}{SL}
\DeclareMathOperator{\PGL}{PGL}
\DeclareMathOperator{\Sym}{Sym}
\DeclareMathOperator{\mult}{mult}
\DeclareMathOperator{\dfc}{def}
\DeclareMathOperator{\per}{per}
\DeclareMathOperator{\rank}{rank}
\DeclareMathOperator{\adj}{adj}
\DeclareMathOperator{\codim}{codim}
\DeclareMathOperator{\Stab}{Stab}
\newcommand{\bC}{\mathbb{C}}
\newcommand{\bP}{\mathbb{P}}
\newcommand{\bZ}{\mathbb{Z}}
\newcommand{\bQ}{\mathbb{Q}}
\newcommand{\cO}{\overline{\mathcal{O}}}
\newcommand{\Rr}{R}
\newcommand{\Pp}{P}
\newcommand{\Dpad}{D^{\mathrm{pad}}}
\newcommand{\Ddet}{D^{\det}}

\title[The determinant's five-row ideal and the padded permanent at $n=4$]
{The onset of the determinant's five-row ideal,\\ and the padded permanent at $n=4$}

\author{Swami Sethuraman}
\address{Beneficus AI}
\email{swsethuraman@beneficus.ai}
\date{\today}
\subjclass[2020]{Primary 14L24; Secondary 68Q17, 13A50, 20G05, 14M12}
\keywords{geometric complexity theory, permanent versus determinant, orbit
closure, determinantal hypersurface, multiplicity obstruction, plethysm}

\begin{document}

\begin{abstract}
Let $\Ddet_r\subset\Sym^{n}\bC^{r}$ be the variety of $r$-variable degree-$n$
forms with an $n\times n$ linear determinantal representation, and let
$\Rr_r=\{\ell\cdot c\}\subset\Sym^{4}\bC^{r}$ be the reducible quartics, the
locus into which the padded permanent $x_0\per_3$ restricts.  Working in the
length-reduced model of geometric complexity theory, we study the ideals of
these varieties as $\GL_r$-modules and the multiplicity gap that a separation
of the padded permanent from the $4\times4$ determinant would require.  We
prove: (i) a \emph{washout theorem} --- the padded permanent restricts to a
generic reducible form at every length $r\le5$, so no covariant of length at
most five can distinguish it from $\ell\cdot(\text{any cubic})$, and the
permanent's own structure first appears at length six and degree seven; (ii) a
\emph{transfer lemma} --- the reducible locus over-estimates the padded
permanent's coordinate ring, so a computation against $\{\ell\cdot c\}$ is a
complete screen for multiplicity obstructions; and (iii) a combinatorial
criterion for membership in the ideal of $\{\ell\cdot c\}$, recovering the
Kadish--Landsberg padding bound as its degenerate case.  We also prove,
modulo three adopted inputs from the literature (Kleiman's transversality
theorem, Dimca's syzygy--defect theorem and the Gulliksen--Neg\aa rd
resolution), (iv) a
\emph{cap theorem} --- for every $n\ge2$ the ideal of $\Ddet_5$ contains the
minors of size $\mathrm{cap}(n)=5n(n-1)^{2}(7n-8)/12$ of the degree-$(3n-5)$
Jacobian (Macaulay) matrix, the rank of that matrix at a smooth form, produced by the $\nu(n)=n^{2}(n^{2}-1)/12$
nodes of the generic determinantal member.  The cap is exact at $n=2$; we
conjecture it is the onset of the ideal for all $n$, and record the
$n=5$ anomaly at which the node count first exceeds the codimension.  A
stabiliser-isotypic reduction makes the underlying multiplicities computable
far beyond a naive sweep.  Across the entire computed range of the padded
$n=4$ comparison we find no multiplicity obstruction, in agreement with, and extending to small parameters,
the no-occurrence-obstruction phenomenon of B\"urgisser--Ikenmeyer--Panova.
\end{abstract}

\maketitle

\section{Introduction}\label{sec:intro}

Valiant's conjecture, the algebraic analogue of $\mathrm{P}\neq\mathrm{NP}$,
asserts that the permanent $\per_m$ cannot be written as the determinant of a
matrix of affine linear forms of size polynomial in $m$
\cite{Valiant,MS,BI,Landsberg}.  In the border, padded formulation one asks
whether the padded permanent $x_0^{\,n-m}\per_m$ lies in the orbit closure
$\overline{\GL_{n^2}\cdot\det_n}$; geometric complexity theory
(Mulmuley--Sohoni \cite{MS}) proposes to obstruct such a containment by a
representation-theoretic inequality.  Containment of orbit closures forces, for
every partition $\lambda$, a comparison of the multiplicities of the
irreducible $\SL$-module $S_\lambda$ in the two coordinate rings; a
\emph{multiplicity obstruction} is a $\lambda$ at which the padded permanent's
multiplicity strictly exceeds the determinant's.  Occurrence obstructions --- the
special case where the determinant's multiplicity is zero --- were shown by
B\"urgisser, Ikenmeyer and Panova \cite{BIP} to be insufficient in the padded,
polynomial-$n$ regime; multiplicity obstructions remain the live hope, and are
strictly stronger \cite{DIP}.

This paper is a study of the smallest genuinely $4\times4$ instance of the
mechanism.  Rather than the full orbit closure we work in the
\emph{length-reduced model}: a highest-weight vector of weight $\lambda$ with
$\ell(\lambda)=r$ rows sees a form only through its restriction to an
$r$-dimensional subspace of the variables, and the multiplicity of $S_\lambda$
in the coordinate ring of an orbit closure equals the multiplicity of
$S_\lambda(\bC^r)$ in the coordinate ring of the restricted variety
\cite[\S4]{Companion}.  The two players become, inside $\Sym^{4}\bC^{r}$,
\[
  \Ddet_r=\overline{\{\det_4(s_1A_1+\dots+s_rA_r):A_i\in M_4\}},
  \qquad
  \Rr_r=\overline{\{\ell\cdot c:\ell\in\bC^r{}^{*},\ c\in\Sym^{3}\bC^{r}\}},
\]
the quinary (etc.) quartics with a $4\times4$ linear determinantal
representation, and the reducible quartics with a linear factor.  The padded
permanent's restrictions land in $\Rr_r$; we write $\Pp_r$ for the closure of
its own restrictions, so $\Pp_r\subseteq\Rr_r$ always.  A multiplicity
obstruction is a weight at which $\mult_\lambda\bC[\Pp_r]>\mult_\lambda\bC[\Ddet_r]$.

\subsection*{Results.}
Our first group of results locates where such an obstruction can live.

\emph{The washout theorem} (\S\ref{sec:washout}).  The restrictions of $\per_3$
to $r$-planes are dense in $\Sym^{3}\bC^{r}$ for every $r\le5$, so
$\Pp_r=\Rr_r$ there and \emph{no covariant of length at most five sees the
permanent}: every separation at $\ell(\lambda)\le5$ is a statement about
reducibility, not about the permanent --- although a positive gap at length five
would still be a multiplicity obstruction for the padded permanent, since
$\Pp_5=\Rr_5$.  The permanent first enters at $r=6$,
where $\dim\Pp_6=55<61=\dim\Rr_6$, and by a Pieri bound only from degree seven
(Proposition~\ref{prop:pieri}).  With the finite-generic-stabiliser count of
\S\ref{sec:washout} the dimension tables of the whole programme become
unconditional.

\emph{The washout persists past the dimension gap.}  At $r\le5$ the equality
$\Pp_r=\Rr_r$ makes the blindness a theorem.  At $r\ge6$ the two varieties
differ, and yet in every weight we have measured the two multiplicities still
agree (measured): $\mult_\lambda\bC[\Pp_r]=\mult_\lambda\bC[\Rr_r]$ at $r=6$ in $682$
six-row cells through degree twelve, and at $r=9$ at the goal cell
$\lambda=(65,17,2^{7})$ and at rung thirteen of its ladder, where the padded and
reducible kernels are the same space over both house primes.  At the goal cell
the \emph{unpadded} $\per_4$ has no equation at all in that weight.  The
mechanism is Proposition~\ref{prop:pieri}'s second half: a strict inequality
$\mult_\lambda\bC[\Pp_r]<\mult_\lambda\bC[\Rr_r]$ in degree $\delta$
requires $I(D^{\per_3}_r)_\delta\neq0$ on the cubic side, where
$D^{\per_3}_r\subseteq\Sym^{3}\bC^{r}$ is the closure of the restrictions of
$\per_3$ to $r$-planes.  That ideal is zero for \emph{every} $r$ in every degree
$\delta\le8$, and at $r=6$ also in degree nine (proved; a reconciliation of
several computations: lengths at most five by Theorem~\ref{thm:washout}, length
six through degree nine, lengths seven and eight at degree eight by exhaustive
exact computation, and the fact that every constituent of $I(D^{\per_3}_r)_\delta$
has length between $6$ and $\min(r,\delta)$).  Hence
$\mult_\lambda\bC[\Pp_r]=\mult_\lambda\bC[\Rr_r]$ at every weight of every
length in every degree at most eight.  So the reducibility reading of the
$r\le5$ washout survives into the range where the varieties are known to be
distinct.

\emph{The transfer lemma} (\S\ref{sec:transfer}).  Because $\Pp_r\subseteq\Rr_r$,
one has $\mult_\lambda\bC[\Pp_r]\le\mult_\lambda\bC[\Rr_r]$ weight by weight; a
multiplicity gap computed against $\Rr_r$ therefore transfers to the padded
permanent when it is negative and must be re-examined when it is positive, and
a computation against the reducible locus is a \emph{complete screen}: no
obstruction hides from it.  The permanent can only \emph{erase} an
$\Rr_r$-obstruction, never create one.

\emph{The reducible ideal} (\S\ref{sec:reducible}).  A highest-weight vector
lies in $I(\Rr_r)$ if and only if each of its coefficient monomials omits some
variable (Theorem~\ref{thm:star}); this makes membership decidable by
inspection, gives a point-free formula for $\mult_\lambda\bC[\Rr_r]$, and
recovers the Kadish--Landsberg padding bound \cite{KL2}: no obstruction can
live at a weight with $\lambda_1<\delta$.  For quartics the ideal
$I(\Rr_r)$ begins in degree five for every $r\ge5$.

Our second group concerns the determinant's own ideal.

\emph{The cap theorem} (\S\ref{sec:cap}; proved modulo Kleiman, Dimca and
Gulliksen--Neg\aa rd, all adopted).  For every $n\ge2$ the ideal of
$\Ddet_5\subset\Sym^{n}\bC^{5}$ contains the size-$\mathrm{cap}(n)$ minors of the
degree-$(3n-5)$ Jacobian (Macaulay) matrix of the five partial derivatives,
where
\[
  \mathrm{cap}(n)=\dim\Sym^{3n-5}\bC^{5}-\mu_{3n-5}(n)=\tfrac{5}{12}n(n-1)^{2}(7n-8)
  =5,\ 65,\ 300,\ 900,\ 2125,\dots
\]
The mechanism is uniform: the generic determinantal member has
$\nu(n)=n^{2}(n^{2}-1)/12$ ordinary double points which fail to impose
independent conditions on forms of degree $2n-5$ by exactly one, and Dimca's
theorem \cite{Dimca13} turns that failure into one extra dimension of the
Milnor algebra in degree $3n-5$.  The cap is exact at $n=2$ (it is the
discriminant of quinary quadrics), is at $n=3$ the degree-$65$ Jacobian cap of
\cite[Prop.~4.23]{Companion}, which improved the discriminant bound $80$ there,
and equals the geometrically motivated $300$ at $n=4$.  For the onset
$\delta_0$ of the length-five part of the ideal of $\overline{\GL_9\cdot\det_3}$
--- the quinary-cubic case --- the current bracket is $6\le\delta_0\le65$
unconditionally and $8\le\delta_0\le65$ given the measured totals of
\cite{Companion}.  We conjecture that the cap is the onset of $I(\Ddet_5)$ at
every $n$ (Conjecture~\ref{conj:onset}); it survives every cell we have
measured.  At $n=5$ the coincidence ``nodes $=$ codimension'' of $n=3,4$ first
fails, and $\Ddet_5$ becomes a superabundant component of the $50$-nodal locus,
by an argument that assumes the minor ideal saturated in degree $n$, which is
measured at $n=3,\dots,7$ (\S\ref{sec:anomaly}).

\emph{The computational engine} (\S\ref{sec:reduction}).  A
stabiliser-isotypic reduction (Theorem~\ref{thm:reduction}) confines the
highest-weight vectors of a weight $\lambda$ to a single character-isotypic
component of the action of the weight-stabiliser, cutting the working
dimension from the weight-space dimension $N_S$ to the dimension $n_\chi$ of
that component (computed exactly, not estimated by $N_S/|\Stab_W(\lambda)|$)
and putting weights of a million-plus monomials within reach of a single
workstation.

\emph{Negative results} (\S\ref{sec:negative}).  On the slab of weights of
length at most four the obstruction gap is non-positive in every degree, by
containment (Theorem~\ref{thm:slab}(1)); the determinant's multiplicity equals
the ambient plethysm bound at length at most three in every degree, and at
length four below the onset $e$ of the principal ideal $I(\Ddet_4)$, with
$e\ge10$ certified and $e=320112$ adopted from \cite{LLV}
(Theorem~\ref{thm:slab}(2)--(3)); the arithmetic route to an obstruction --- occurrence
cells, where the determinant's coordinate ring omits a weight, and forced
cells, where the ambient bound already exceeds the stabiliser bound --- is silent
across the entire computed range at $n=4$, extending to small parameters the
no-occurrence phenomenon of \cite{BIP}; and no multiplicity obstruction was
found in any measured cell.  The set-theoretic non-containment
$x_0\per_3\notin\overline{\GL_{16}\cdot\det_4}$ holds independently, by
\cite[Thm.~1.0.1]{LMR} ($\overline{dc}(\per_3)\ge9/2$), so, as in \cite{DIP}, the programme exhibits the
multiplicity method on a separation that is known by other means; the value is
in charting where the method can and cannot certify it.

Throughout, we distinguish sharply between statements that are \emph{proved},
those \emph{adopted} from the literature (Kleiman's transversality theorem,
Dimca's syzygy theorem, the Gulliksen--Neg\aa rd resolution, and the degree
$e=320112$ of the determinantal quartic surfaces \cite{LLV}), and those
\emph{measured} by exact computation over two word-size primes with
$\rank(R)=\dim-\text{(kernel)}$ asserted as an internal check.  The
conjectures are labelled as such.  For the external results used in proofs we
record how each was checked: read in the primary source ---
\cite[(1.9), Cor.~6.4]{Beau}, \cite[Thm.~1.3]{KL2}, \cite[Thms.~1.0.1, 2.3.1,
3.1.1 and \S3.2]{LMR}, \cite[Thm.~2, Lemma~4.2, Table~2, Cor.~3.1]{LLV}, and
\cite[Thm.~3.1]{Dimca13} at the level of its statement; adopted from secondary
accounts without a primary reading --- \cite{Kleiman} and \cite{GN}; and not
read --- \cite{BrunsVetter}, cited with \cite{GN} for the same resolution.  All
code,
inputs and exact outputs are available (\S\ref{sec:data}).

\section{The two varieties and the length reduction}\label{sec:varieties}

Fix $n=4$ unless stated otherwise and write $W_r=\Sym^{4}\bC^{r}$.  A weight
$\lambda$ with $\ell(\lambda)=r$ contributes to the coordinate ring of an
orbit closure $\cO\subset\Sym^{4}\bC^{16}$ exactly through the restriction to
an $r$-plane, and the length reduction (\cite[Prop.~4.19]{Companion} for
$\det_3$; the proof uses only that $\lambda$ has $r$ rows and applies verbatim
to $\det_4$ and to $x_0\per_3$) identifies
\begin{equation}\label{eq:lengthred}
  \mult_\lambda\bC[\cO]_\delta=\mult_{S_\lambda(\bC^r)}\bC[\cO|_{\bC^r}]_\delta,
\end{equation}
where $\cO|_{\bC^r}$ is the restricted variety.  For $\cO=\overline{\GL_{16}\cdot\det_4}$
this is $\Ddet_r$; for the padded permanent $\cO=\overline{\GL_{16}\cdot(x_0\per_3)}$
it is $\Pp_r$.  The reducible locus $\Rr_r$ is the ambient container of $\Pp_r$.

\begin{proposition}[Dimensions]\label{prop:dims}
For $r\ge3$, $\dim\Ddet_r\le\min(\dim W_r,\,16r-30)$, with equality for
$3\le r\le6$ (Lemma~\ref{lem:finstab} and the table in \S\ref{sec:washout});
concretely $\dim\Ddet_4=34$ (a hypersurface in $W_4$, $\dim W_4=35$; compare
\cite[Cor.~3.1]{LLV}: linear determinantal surfaces of degree $d$ form a family
of projective dimension $2d^{2}+1$), $\dim\Ddet_5=50$ and $\codim\Ddet_5=20$ in
$W_5$ $(\dim W_5=70)$; $\dim\Ddet_6=66$.  For the padded permanent, $\dim\Pp_r=r+\dim\Sym^{3}\bC^{r}-1$
for $r\le5$ (Corollary~\ref{cor:washout}) and $\dim\Pp_6=55$, while
$\dim\Rr_6=61$; $\dim\Rr_r=r+\binom{r+2}{3}-1$.
\end{proposition}

The proofs are Jacobian rank computations at explicit points together with the
finite-stabiliser upper bounds of \S\ref{sec:washout}; the tables were verified
at both primes and are unconditional given that section.  The generic member of
$\Ddet_5$ is a quartic threefold with exactly $20$ ordinary double points and
defect one (\S\ref{sec:cap}; the node count uses Kleiman's theorem, adopted);
it is non-$\bQ$-factorial, the $n=4$ analogue of
the six-nodal cubic threefold of \cite{Companion}.

\section{The washout theorem}\label{sec:washout}

The following is elementary but decisive: it explains why the length-five hunt,
where computation is cheapest, cannot see the permanent at all.

\begin{theorem}[Washout]\label{thm:washout}
For $r\le5$ the restriction map $\GL_r\cdot\per_3\to\Sym^{3}\bC^{r}$ is
dominant; equivalently $\Pp_r=\Rr_r=\{\ell\cdot c:c\in\Sym^{3}\bC^{r}\}$.
Consequently, for every weight $\lambda$ with $\ell(\lambda)\le5$ and every
$\delta$,
\[
  \mult_\lambda\bC[\Pp_r]=\mult_\lambda\bC[\Rr_r],
\]
and no $\GL_r$-covariant of length at most five distinguishes the padded
permanent from $\ell\cdot(\text{an arbitrary cubic})$.
\end{theorem}

\begin{proof}
A single point of $\Sym^{3}\bC^{5}$ at which the differential of
$g\mapsto\per_3(g\cdot)|_{\bC^{5}}$ has full rank $35$ proves dominance, since
the image is then dense; such a point is exhibited (rank $35$ at an explicit
integer substitution, both primes).  For $r<5$ restrict further.  As
$\per_3$ is a cubic, its restrictions are cubics, and dominance makes them fill
$\Sym^{3}\bC^{r}$; multiplying by the padding linear form gives all of
$\Rr_r$.
\end{proof}

\begin{corollary}\label{cor:washout}
The onset of the determinant's separation from the permanent is a
\emph{length-six} phenomenon.  At $r=6$ the restrictions of $\per_3$ are no
longer dense: $\dim\Pp_6=6+\dim D^{\per_3}_6-1=6+50-1=55$, whereas
$\dim\Rr_6=6+56-1=61$.
\end{corollary}

The upper bound $\dim D^{\per_3}_6\le50$ requires a four-dimensional group
preserving $\per_3$ under which a generic $6$-tuple of restrictions has finite
stabiliser; the lower bound is a Jacobian rank of $50$ at an explicit point.
We record the stabiliser fact once, as it underlies every dimension table.

\begin{lemma}[Finite generic stabiliser]\label{lem:finstab}
Let $G=\{(P,Q)\in\GL_4\times\GL_4:\det P\det Q=1\}$ act on $r$-tuples of
$4\times4$ matrices by $A_i\mapsto PA_iQ$, and let $T$ be the pairs of diagonal
$3\times3$ matrices $(D_1,D_2)$ with $\det D_1\det D_2=1$ acting on $r$-tuples
of $3\times3$ matrices by $A_i\mapsto D_1A_iD_2$.  Both preserve the relevant form
($\det_4$, resp.\ $\per_3$, since $\per(D_1AD_2)=\det D_1\det D_2\per A$),
their effective quotients by the scalars $(\mu I,\mu^{-1}I)$ have dimensions
$30$ and $4$, and for $r\ge3$ (resp.\ every $r\ge1$) the generic $r$-tuple has
trivial stabiliser in the effective group.  Hence
\[
  \dim\Ddet_r\le\min\!\big(\dim W_r,\ 16r-30\big)\quad(r\ge3),\qquad
  \dim D^{\per_3}_r\le\min\!\big(\tbinom{r+2}{3},\ 9r-4\big)\quad(r\ge1),
\]
where $D^{\per_3}_r\subseteq\Sym^{3}\bC^{r}$ is the closure of the restrictions
of $\per_3$ to $r$-planes.  The case $r=2$ is exceptional for $\det_4$.
\end{lemma}

\begin{proof}
If a group $H$ acts on the parameter space of a map constant on $H$-orbits and
the generic stabiliser is finite, the image closure has dimension at most the
parameter dimension minus $\dim H$, since the generic fibre contains an orbit
of dimension $\dim H$.  For $\det_4$ with $r\ge3$: normalise $A_1=I$; then
$PA_1Q=A_1$ forces $Q=P^{-1}$, and $P$ commutes with $A_2,\dots,A_r$.  Two
generic $4\times4$ matrices generate $M_4$ as an algebra (take $A_2$ with
distinct eigenvalues, whose polynomials give the diagonal idempotents $E_{ii}$
in its eigenbasis, and $A_3$ with all entries nonzero in that basis, so that
$E_{ii}A_3E_{jj}$ is a nonzero multiple of $E_{ij}$), so $P$ is scalar.  For
$\per_3$: if $D_1A_1D_2=A_1$ with every entry of $A_1$ nonzero, then
$d_je_k=1$ for all $j,k$, so $(D_1,D_2)$ is scalar.  At $r=2$, after
normalising $A_1=I$, the stabiliser is the commutant of the single generic
matrix $A_2$, its $4$-dimensional algebra of polynomials, $3$-dimensional
modulo scalars; the bound would read $\dim\Ddet_2\le2$ and is false
($\Ddet_2=\Sym^{4}\bC^{2}$, of dimension $5$).  Only these tori and groups are
used; the full stabilisers of $\det_4$ and $\per_3$ are not needed.  The
$n=3$ determinant case is in \cite[\S4]{Companion}.
\end{proof}

The groups in Lemma~\ref{lem:finstab} have the dimensions of the full
stabilisers of $\det_4$ and $\per_3$, but the lemma gives only upper bounds on
dimension; each entry of the table below is an equality because a Jacobian rank
at an explicit point meets the bound.  For other $n$ the analogous count
$\dim\Ddet_5=3n^{2}+2$ is carried as adopted beyond $n=4$ and is used only in
\S\ref{sec:anomaly}.

With Lemma~\ref{lem:finstab} the codimension table
\[
\begin{array}{c|ccc}
 r & \dim W_r & \dim\Ddet_r & \dim\Rr_r\\\hline
 3 & 15 & 15 & 12\\
 4 & 35 & 34 & 23\\
 5 & 70 & 50 & 39\\
 6 & 126& 66 & 61
\end{array}
\]
is unconditional; the pad row is exact because $\per_3$ fills all of
$\Sym^{3}\bC^{r}$ for $r\le5$, and the determinant row is pinned by a sandwich
of a measured Jacobian rank (a lower bound) against the stabiliser count (an
upper bound), which meet.

\begin{proposition}[Pieri delay]\label{prop:pieri}
$I(\Pp_6)_\delta=I(\Rr_6)_\delta$ for $\delta\le5$; the permanent's own
structure is first visible in degree $\ge7$.  More precisely, an equation of
$\Pp_r$ of degree $\delta$ that is not an equation of $\Rr_r$ pulls back to an
equation of $D^{\per_3}_r\subset\Sym^{3}\bC^{r}$ of degree
$\delta$; for $r=6$ this ideal is empty in degree $\le5$ (an equation of degree
$\delta$ has at most $\delta$ rows, and $D^{\per_3}_5=\Sym^{3}\bC^{5}$
by Theorem~\ref{thm:washout}), and is measured empty in degree $6$.
\end{proposition}

\section{The transfer lemma}\label{sec:transfer}

\begin{theorem}[Transfer]\label{thm:transfer}
Set $D_\lambda(X)=\mult_\lambda\bC[X]$.  For every $r$ and $\lambda$,
$D_\lambda(\Pp_r)\le D_\lambda(\Rr_r)$.  Consequently, writing
$\Delta(X)=D_\lambda(\Pp_r)-D_\lambda(\Ddet_r)$ against a fixed determinant
side:
\begin{enumerate}
\item a cell with $D_\lambda(\Rr_r)-D_\lambda(\Ddet_r)<0$ (an
``$\Rr$-blindness'' cell) has $\Delta<0$ as well;
\item a cell with $D_\lambda(\Rr_r)-D_\lambda(\Ddet_r)>0$ need \emph{not}
transfer to $\Pp_r$ and must be re-evaluated against $\Pp_r$;
\item a computation against $\Rr_r$ is a complete screen for obstructions:
if $\Delta>0$ then $D_\lambda(\Rr_r)-D_\lambda(\Ddet_r)>0$.
\end{enumerate}
The evaluation pipeline, which tests at true padded-permanent points, computes
$\Pp_r$ directly; the caveat bites only on literature-style arguments about the
reducible locus.
\end{theorem}

\begin{proof}
$\Pp_r\subseteq\Rr_r$ gives $I(\Rr_r)\subseteq I(\Pp_r)$, hence the ideal
multiplicities satisfy the reverse inequality and the coordinate-ring
multiplicities $D_\lambda(\Pp_r)\le D_\lambda(\Rr_r)$ term by term.  Items
(1)--(3) are immediate, (3) being the contrapositive of (1).
\end{proof}

The lemma is what makes the reducible-locus analysis of the next section a tool
against the permanent rather than a digression: every obstruction is an
$\Rr_r$-obstruction, and the permanent's role is confined to possibly erasing
some of them.

\section{The reducible ideal and the padding bound}\label{sec:reducible}

Let $c_\alpha$ denote the coefficient functional extracting the $x^\alpha$
coefficient, $|\alpha|=n$.  A monomial in the $c_\alpha$ is said to
\emph{omit} the variable $x_i$ if some factor $c_\alpha$ has $\alpha_i=0$.

\begin{theorem}[The reducible criterion]\label{thm:star}
Let $v\in\bC[\Sym^{n}\bC^{r}]_\delta$ be a highest-weight vector of weight
$\lambda$.  Then $v$ vanishes on $\Rr_r=\{\ell\cdot c\}$ if and only if every
monomial of $v$ omits $x_i$ for every $i$.  More generally, for the locus
$\{\ell^{k}\cdot c\}$ of forms with a linear factor of multiplicity $k$, the
criterion reads: every monomial has, for each $i$, a factor $c_\alpha$ with
$\alpha_i<k$.
\end{theorem}

\begin{proof}[Proof sketch]
By the Bruhat decomposition it suffices to test $v$ on the dense cell
$G\cdot(x_1\cdot c)$, $c$ arbitrary; since $v$ is a $B$-eigenvector, this
reduces to vanishing on $x_i\cdot(\text{all cubics})$ for each $i$, i.e.\ on
the coordinate subspace of forms divisible by $x_i$.  A monomial in the
$c_\alpha$ vanishes on that subspace precisely when each factor sees the
divisibility, which is the stated omission condition.  The multiplicity-$k$
version replaces $\alpha_i=0$ by $\alpha_i<k$.
\end{proof}

\begin{corollary}[Point-free multiplicity]\label{cor:pointfree}
$\mult_\lambda\bC[\Rr_r]_\delta=a(\lambda,\delta)-\dim\{v:\text{(\ref{thm:star})}\}$,
where $a(\lambda,\delta)$ is the ambient plethysm multiplicity of $S_\lambda$
in $\bC[\Sym^{n}\bC^{r}]_\delta$.  In particular the reducible-locus side of
every cell is computable by linear algebra on monomial supports, with no
evaluation points.
\end{corollary}

\begin{corollary}[Padding bound]\label{cor:kl}
If $\lambda_1<\delta$ then every monomial of every weight-$\lambda$
highest-weight vector omits every variable $x_i$ automatically (its $\delta$
factors contribute $\lambda_i\le\lambda_1<\delta$ to the $i$-th coordinate, so
some factor has $\alpha_i=0$), so every such vector vanishes on $\Rr_r$:
$\mult_\lambda\bC[\Rr_r]_\delta=0$ and $I(\Rr_r)_\delta$ contains the whole
$\lambda$-isotypic component.  Hence $\mult_\lambda\bC[\Pp_r]_\delta=0$ and
\emph{no multiplicity obstruction can live at a weight with $\lambda_1<\delta$}.
\end{corollary}

Corollary~\ref{cor:kl} is the first assertion of \cite[Thm.~1.3]{KL2} in the
case $n-m=1$: the ideal of the padded polynomials $\ell^{\,n-m}h$ contains the
whole isotypic component of $S_\lambda$ in degree $\delta$ whenever
$\lambda_1<\delta(n-m)$.  It is recovered here from the combinatorial
criterion; the technique --- evaluating highest-weight vectors at
$x_1^{\,n-m}h$ and reading off their monomial support --- is theirs.  In the
sources checked in one literature pass, among them \cite{KL2} and
\cite{Landsberg}, we did not find the exact criterion of
Theorem~\ref{thm:star} stated; we record it as not found there, not as new.  For quartics, the ideal
$I(\Rr_r)$ begins in degree five for every $r\ge5$: the unique degree-five
invariant of quinary quartics vanishes on $\{\ell\cdot c\}$ (it is an $a=1$
cell, invisible to the $a\ge2$ gate discussed in \S\ref{sec:negative}), and no
covariant of lower degree does.

\section{Containment at $r\le4$, non-containment at $r=5$}\label{sec:contain}

The obstruction-relevant comparison changes character at length five.

\begin{proposition}[Containment, $r\le4$]\label{prop:contain}
For $r\le4$, $\Rr_r\subseteq\Ddet_r$, hence
$\mult_\lambda\bC[\Rr_r]_\delta\le\mult_\lambda\bC[\Ddet_r]_\delta$ and
$\Delta\le0$ at every weight of length at most four and every degree.
\end{proposition}

\begin{proof}
A general quaternary cubic is $3\times3$ linear determinantal: every smooth
cubic surface is (\cite[Cor.~6.4]{Beau}, classically Grassmann's), and
independently the restriction map $M\mapsto\det_3(s_1A_1+\dots+s_4A_4)$ has
full Jacobian rank $20$ at an explicit integer point, so its image is dense in
$\Sym^{3}\bC^{4}$.  For such a cubic $c=\det_3M(s)$,
$\ell\cdot c=\det_4\operatorname{diag}(\ell,M)$ is a $4\times4$ linear
determinantal representation.  $\Ddet_4$ is closed, so it contains the closure
of these reducible quartics, which is $\Rr_4$; for $r\le3$ restrict further
(or note $\Ddet_r=\Sym^{4}\bC^{r}$).  Containment gives a surjection
$\bC[\Ddet_r]\twoheadrightarrow\bC[\Rr_r]$ of $\GL_r$-algebras, hence the
inequality of multiplicities, and $\Pp_r\subseteq\Rr_r$ gives $\Delta\le0$.
The argument does not extend to $r=5$: a generic cubic threefold has no
non-trivial determinantal representation at all, since by the
Noether--Lefschetz argument of \cite[(1.9)]{Beau} a generic hypersurface of
degree $d$ in $\bP^{m}$ is a non-trivial determinant only if $m=2$, or $m=3$
and $d\le3$; in the present setting the image of $3\times3$ linear
determinants in $\Sym^{3}\bC^{5}$ has dimension $29<35$ (Jacobian rank, both
primes).
\end{proof}

Write $W=s_5\cdot\Sym^{3}\bC^{5}\subset W_5$ for the $35$-dimensional space of
quartics divisible by a fixed linear form; since every element of $\Rr_5$ is
$\GL_5$-equivalent to an element of $W$, $\Rr_5\subseteq\Ddet_5$ if and only if
$W\subseteq\Ddet_5$.

\begin{theorem}[Non-containment at $r=5$, determinant part]\label{thm:noncontain}
If $\det(s_1A_1+\dots+s_5A_5)\in W$ for a pencil of $4\times4$ matrices, then
it equals $s_5\cdot C$ with $C$ either a $3\times3$ linear determinantal cubic
or a cubic containing a plane.  These two families of cubics have dimensions
$29$ and $31$, so the actual determinantal quartics in $W$ form a set of
dimension at most $31<35=\dim W$.
\end{theorem}

Theorem~\ref{thm:noncontain} rests on a classification of the singular
four-dimensional subspaces of $M_4$: if $s_5$ divides $\det M(s)$ then
$\operatorname{span}(A_1,\dots,A_4)$ consists of singular matrices.  Those
subspaces are not all compression spaces --- that is the obstacle the
classification had to overcome, not the reason for the conclusion --- and the
exceptional strata give cubic families of dimensions $27$ and $25$, below the
compression maximum $31$ \cite{Companion2}.  The separation is the dimension
count $31<35$; no highest-weight vector enters.

The theorem is a statement about actual determinants.  It does not by itself
give $\Rr_5\not\subseteq\Ddet_5$, because $\Ddet_5$ is a closure: a point of
$\Ddet_5\cap W$ may be a limit of determinants that is not itself a
determinant, and no classification of such boundary points is available here
(Remark~\ref{rem:noncontain-status}).  The closure statement
$\Rr_5\not\subseteq\Ddet_5$ is carried in this paper as \emph{adopted}: it is a
separate, unpublished result of this programme, carried as ADOPTED on the programme's
record, that exhibits a padded smooth cubic threefold
$s_5\cdot C^{\ast}$ outside $\Ddet_5$ \cite{Companion2}; it is not a
consequence of Theorem~\ref{thm:noncontain}.  Given it, together with the
dimension gap $39<50$ of the codimension table, the two varieties differ in
both directions at $r=5$ --- yet, as \S\ref{sec:negative} records, neither
difference has yet been made to produce a weight with $\Delta>0$.

\begin{remark}[What the non-containment currently rests on]\label{rem:noncontain-status}
Two different objects are enumerated in this section, and they should not be
confused.  Theorem~\ref{thm:noncontain} classifies the singular
four-dimensional subspaces of $M_4$, which governs the \emph{actual}
determinants in $W$.  This remark concerns an attempt to reach the closure
through the blow-up of the base scheme of the parametrisation $\Phi$, whose
exceptional divisor is governed by the components of
$\operatorname{Proj}\operatorname{gr}_J R$; a classification of singular
subspaces does not by itself classify those components.  The computational
record has sharpened both what is established and what is not.  Every measured exceptional image has since been pushed below the
threshold: through the reachable contact orders, and at every incidence of the
primitive family, no reducible image exceeds $29<31\ll35$.  The following are
now proved rather than measured.

\emph{(i) The contact-order lemma.}  At a smooth point of $V(J)$ the
exceptional fibre is $\bP(\operatorname{im} d\Phi)$ at every order.  This turns
the observed invariance of contact order into a theorem and, more usefully,
localises any hidden component of the exceptional divisor to $\mathrm{Sing}\,V(J)$.

\emph{(ii) The stratum list of the base locus of $\Phi$ is longer than first
enumerated.}  (This is a list of components of a base locus, not of singular
subspaces, and it does not contradict Theorem~\ref{thm:noncontain}.)  The
base locus decomposes into eight components once the singular-pencil stratum
$SP$ is included, and a further $49$-dimensional semi-primitive component was
found that is absent from every earlier stratum list; moreover $\Pp\cap\ker$ is
a rank-$\le2$ locus rather than the rank-$3$ incidence it had been taken for.

\emph{(iii) The base scheme is non-reduced.}  At a generic point of every
pairwise incidence the initial ideal is strictly smaller than that of the
intersection of the component ideals --- $\dim Q_2=12<16$, $41<69$, $25<49$ and
$59<144$ at the four incidences --- so the base scheme of $\Phi$ at $r=5$ is
\emph{generically reduced along every component and non-reduced along every
pairwise incidence}.  This is the precise sense in which the normal cone of $J$
differs from the normal cone of the reduced base locus, and it is why the
exceptional fibre over an incidence is larger than $\bP(\operatorname{im}d\Phi)$.

Consequence (iii) is why the bound must be quantified over the normal cone and
not over the reduced singular locus: a component of the exceptional divisor can
be supported over a proper sublocus of an irreducible singular component, or
over embedded structure, and be invisible at the generic point of any of them.
The remaining obligation is therefore

\begin{quote}
for every irreducible component $E$ of $\operatorname{Proj}\operatorname{gr}_J R$
supported over $\mathrm{Sing}\,V(J)$ --- including components over proper
subloci, over the incidence and rank-degeneration strata, and components arising
from embedded structure --- the fixed-factor image of $E$ has dimension $<35$,
\end{quote}

which together with (i) exhausts the blow-up.

\emph{(iv) The four residues are now numbers, and the interior is exact.}  The
four loci this remark first named as unbounded have since been measured.  The
order-$2$ image over $\Pp\cap c_{21}$ --- the one number the earlier survey left
open --- is $19$.  The rank strata of $M(a)$ at $\ker\cap\operatorname{coker}$
are $0$ and $9$ only, both of image $29$.  The deeper rank-$\le2$ incidences
drop ($28\to19\to9$), and every incidence meeting an exact-type locus inherits
its determinant, so is bounded by $31$.  Contact order $\ge4$ at the incidences
is obstructed on the linear route and the image is saturated at order $3$; that
last one is evidence rather than proof, and is the only one of the four not
closed by an exact computation.

More is true, and it is the half with no closure gap.  Write
$W=s_5\Sym^{3}\bC^{5}$ for the fixed-factor space, $\dim W=35$.  The interior
$W_{\mathrm{int}}=\Phi(X_5)\cap W$ reparametrises through the $r=4$ base locus
$B_4$ --- $s_5\mid\det M(s)$ exactly when $(A_1,\dots,A_4)\in B_4$, with $A_5$
free --- so it is a finite union of images of irreducible parametrised families,
and its dimension is the generic Jacobian rank over each: an \emph{exact upper}
bound, not an arc lower bound.  It is
\[
  \dim W_{\mathrm{int}} = 31 .
\]
This interior count is the content of Theorem~\ref{thm:noncontain}: for the
fixed factor, the cubics $C$ with $s_5C$ an actual determinant are the
$3\times3$ determinantal cubics ($29$) and the cubics containing a plane
($31$) \cite{Companion3}.  With every enumerated exceptional component at most $31$ as well,
$\dim(\Ddet_5\cap W)\le31<35=\dim W$ holds on the enumerated cone.

\emph{What is not proved} is that the enumeration is complete --- that
$\operatorname{Proj}\operatorname{gr}_J R$ has no component outside the list ---
nor, what would settle the closure question directly, a description of the
boundary
$(\Ddet_5\setminus D^{\det,\circ}_5)\cap\Rr_5$, where $D^{\det,\circ}_5$ is the set
of actual determinants.  Any component of that boundary intersection has
dimension at least $50+39-70=19$; nothing here excludes one \cite{Companion3}.  So the bound on
$\dim(\Ddet_5\cap W)$ is established for the determinant part only, and the
closure statement $\Rr_5\not\subseteq\Ddet_5$ is used in this paper only as
the adopted statement recorded after Theorem~\ref{thm:noncontain}.
\end{remark}

\section{The cap theorem}\label{sec:cap}

We now turn to the determinant's own ideal, the object whose onset bounds the
obstruction hunt from below.  Fix $\Ddet_5\subset\Sym^{n}\bC^{5}$ and let
$M_k(F)$ be the degree-$k$ Macaulay matrix of the five partials
$\partial_1F,\dots,\partial_5F$: rows indexed by $(i,m)$ with $m$ a monomial of
degree $k-n+1$, columns by monomials of degree $k$, the entry the coefficient
of the column monomial in $m\,\partial_iF$.  Its entries are linear in $F$, its
row space is $(J_F)_k$, and $\codim(J_F)_k=\dim(\bC[s]/J_F)_k$ is the corank.
For smooth $F$ the partials are a regular sequence and
$\mu_k(n):=\dim(\bC[s]/J_F)_k=[t^k]\big((1-t^{n-1})/(1-t)\big)^{5}$.  Put
\[
  \mathrm{cap}(n)=\dim\Sym^{3n-5}\bC^{5}-\mu_{3n-5}(n)
  =5\binom{2n}{4}-10\binom{n+1}{4}=\tfrac{5}{12}n(n-1)^{2}(7n-8).
\]

\begin{theorem}[The cap; proved modulo Kleiman \cite{Kleiman} (read-status
SECONDARY), Dimca {\cite[Thm.~3.1]{Dimca13}} (PRIMARY, statement level) and
Gulliksen--Neg\aa rd \cite{GN} (SECONDARY), all adopted]\label{thm:cap}
For every $n\ge2$ the size-$\mathrm{cap}(n)$ minors of $M_{3n-5}(F)$ are
nonzero polynomials of degree $\mathrm{cap}(n)$ that vanish identically on
$\Ddet_5$; their span is a nonzero $\GL_5$-submodule of $I(\Ddet_5)_{\mathrm{cap}(n)}$.
Hence the onset of $I(\Ddet_5)$ is at most $\mathrm{cap}(n)$, with values
$5,65,300,900,2125,4305,\dots$ at $n=2,3,4,5,6,7,\dots$.
\end{theorem}

\begin{proof}[Proof]
The case $n=2$ is exact: $M_1(F)$ is twice the symmetric $5\times5$ matrix of
the quadric, $\mathrm{cap}(2)=5$, and the unique $5\times5$ minor is the
discriminant, whose vanishing cuts out the rank-$\le4$ quadrics $=\Ddet_5$; the
ideal is principal in the UFD $\bC[\Sym^{2}\bC^{5}]$, so the onset is exactly
$5$.  Assume $n\ge3$; the argument is four steps.

\emph{Step 1 (every member is singular along the rank locus).} Jacobi's formula
gives $\partial_kF=\operatorname{tr}(\adj M(s)\cdot A_k)$; at a point of the
pencil where $\rank M(s)\le n-2$ the adjugate vanishes, so all partials vanish.

\emph{Step 2 (the generic member has $\nu(n)$ nodes).} By Kleiman's
transversality theorem \cite{Kleiman} (characteristic zero) a generic
$\bP^{4}$ misses the rank-$\le n-3$ locus, meets the rank-$\le n-2$ locus
transversally in $\nu(n)=\deg\{\rank\le n-2\}=n^{2}(n^{2}-1)/12$ reduced
points, and meets the smooth determinant locus transversally elsewhere; a
Schur-complement expansion at such a point exhibits a rank-four quadratic
leading term, i.e.\ an ordinary double point.  So $\operatorname{Sing}F$ is
exactly these $\nu(n)$ nodes.

\emph{Step 3 (the nodes fail forms of degree $2n-5$ by one).} The ideal $J$ of
the submaximal minors of the generic $n\times n$ matrix has the
Gulliksen--Neg\aa rd resolution \cite{GN,BrunsVetter}, which specialises to the
pencil (grade four, generic perfection) and yields the Hilbert function
\[
  H_{S/J}(k)=[t^k]\frac{1-n^{2}t^{n-1}+(2n^{2}-2)t^{n}-n^{2}t^{n+1}+t^{2n}}{(1-t)^{5}}.
\]
A direct evaluation, a polynomial identity in $n$ verified symbolically and
numerically for $3\le n\le10$, gives $H_{S/J}(2n-5)=\nu(n)-1$.  Since
$J\subseteq I_N$ for the node scheme $N$, the failure
$\dfc_{2n-5}(N)\ge\nu(n)-H_{S/J}(2n-5)=1$.  At $n=3$ this is elementary --- six
points impose at most five conditions on the five-dimensional space of linear
forms --- and at $n=4$ it is the sixteen independent cubics through twenty
nodes.

\emph{Step 4 (Dimca).} For a hypersurface of degree $d$ in $\bP^{4}$ with
isolated singularities, Dimca's theorem \cite[Thm.~3.1]{Dimca13} identifies the
degree-$(3d-5)$ piece of the Milnor algebra as $\mu_{3d-5}+\dfc_{2d-5}(\Sigma)$.
With Step 3, $\dim(\bC[s]/J_F)_{3n-5}\ge\mu_{3n-5}(n)+1$, i.e.\
$\rank M_{3n-5}(F)\le\mathrm{cap}(n)-1$ on the dense determinantal locus, hence
on $\Ddet_5$; the minors vanish there.  For smooth $F$ the rank is exactly
$\mathrm{cap}(n)$, so some minor is a nonzero polynomial; equivariance of the
Fitting construction makes the span $\GL_5$-stable, and \eqref{eq:lengthred}
carries it into $I(\cO)$.
\end{proof}

\emph{Measured, both primes, fresh pencils.} The corank of $M_{3n-5}$ on
$\Ddet_5$ is $6,31,102,256,541$ at $n=3,4,5,6,7$ against the smooth
$5,30,101,255,540$; so $\dfc_{2n-5}=1$ exactly and $\mathrm{cap}(n)$ is the
\emph{exact} rank there.  We have not found the values at $n=5,6,7$ computed
elsewhere; we record that, and claim no priority for them.

The cap minors are equations of $\Ddet_5$, but they do not cut it out.  At
$n=3$ the size-$65$ minors of $M_4(C)$ also vanish at every quinary cubic $C$
containing a plane: writing $C=x_1q_1+x_2q_2$, the $2\times2$ minors of
$(\partial_jq_1,\partial_jq_2)_{j=3,4,5}$ give a non-Koszul syzygy among the
partials with quadric coefficients, so $\rank M_4(C)\le64$ on a dense subset of
that family and hence on all of it \cite{Companion3}.  That family has dimension $31$, larger
than $\dim\Ddet_5=29$, so the zero locus of the cap minors strictly contains
$\Ddet_5$: the ideal the minors generate does not cut out $\Ddet_5$, even
set-theoretically, so they do not by themselves identify $I(\Ddet_5)$.  This
does not affect Theorem~\ref{thm:cap} or
Conjecture~\ref{conj:onset}, which concern the onset, but the minors cannot be
used as a separating equation against anything meeting that family.

\begin{conjecture}[The onset is the cap]\label{conj:onset}
For every $n\ge2$, the onset of $I(\Ddet_5)$ equals $\mathrm{cap}(n)$: the
Jacobian minors are the first equations.  Proved at $n=2$; open at $n\ge3$;
falsified by any length-five component in degree below $65$ ($n=3$) or $300$
($n=4$).
\end{conjecture}

The conjecture survives every measurement made: at $n=3$ the ideal is empty
through degree five unconditionally and through degree seven given the
measured totals of \cite{Companion}, and empty on every one of $121$
length-five cells measured at degrees eight and nine (the affordable cells, not
all cells), and the unique
degree-ten $\SL_5$-invariant of cubic threefolds does \emph{not} vanish on
$\Ddet_5$, so no invariant of degree below fifteen is an equation; at $n=4$ the
ideal is empty through degree seven and on every measured cell at degree eight.
The support is structural rather than probabilistic: within the Jacobian family
$3n-5$ is the cheapest degree, the discriminant route costs more at every $n$,
and no equation is forced by any Hilbert-function comparison below the cap.

\subsection{The $n=5$ anomaly}\label{sec:anomaly}

At $n=3,4$ the codimension of $\Ddet_5$ equals the node count $\nu(n)$
($6=6$, $20=20$).  This is the identity $\codim\Ddet_5=\nu(n)-\binom{n-1}{4}$
with $\binom{n-1}{4}=0$ for $n\le4$.  From $n=5$ it fails:
$\nu(5)=50$ but $\codim\Ddet_5=49$, and $\Ddet_5$ is a superabundant component
of the $50$-nodal locus of quintic threefolds, of dimension $76$ against the
expected $75$, the excess being the one condition the fifty nodes fail to
impose on quintics ($2n-5=n$ at $n=5$).  The component statement is proved at
every $n$ at which the ideal of submaximal minors is saturated in degree $n$
(a tangent-space computation on the incidence variety); saturation is measured
at $n=3,\dots,7$ and is an expectation beyond.  The codimension at $n=5$ also
uses the finite-stabiliser count $\dim\Ddet_5=3n^{2}+2$, adopted beyond
$n=4$.  The cap theorem is untouched: it uses
only $\dfc_{2n-5}\ge1$, and a larger defect only lowers the cap.

\subsection{The $n=3$ twin}\label{sec:twin}

The $n=3$ instance of the mechanism is sharp, and answers, for nodes in general
position, the first part of \cite[Question~8.5]{Companion}, where the same
argument is also sketched.  (Beauville's treatment of determinantal
hypersurfaces \cite{Beau} restricts to smooth hypersurfaces and does not
contain it; a wider search of the classical literature on nodal cubic
threefolds has not been made, and we claim no priority.)

\begin{theorem}\label{thm:frame}
$\Ddet_5\subset\Sym^{3}\bC^{5}$ is the closure of the cubic threefolds singular
at six points in linearly general position.  Equivalently, the generic
six-nodal cubic threefold with nodes in general position is determinantal, and
$\Ddet_5$ is an irreducible component of the closure of the six-nodal locus.
\end{theorem}

\begin{proof}[Proof sketch]
Six points in linearly general position form a projective frame, unique up to
$\PGL_5$; the cubics singular at a fixed frame form a $\bP^{4}$ (the $30$ node
conditions are independent, checked exactly), so the six-nodal-in-general-position
locus is the irreducible image of $\PGL_5\times\bP^{4}$, of dimension
$24+4=28=\dim\bP(\Ddet_5)$.  It remains to show the six nodes of a generic
determinantal cubic are a frame; the frame condition is open on the Grassmannian
of pencils, and holds at one explicit pencil (whose sixth node is computed
exactly by a Gr\"obner basis and verified to make every five of the six span
$\bP^{4}$).  Both varieties being irreducible closed of dimension $28$, they
coincide.  For the component statement: at a generic determinantal cubic the
six nodes impose independent conditions on cubics ($\dfc_3(N)=0$), so the
incidence variety of cubics with six marked singular points is smooth there
with tangent space of dimension $28$, and $\Ddet_5$ is the unique component
through that point (\S\ref{sec:anomaly}, the saturated case $n=3$).
\end{proof}

\section{The stabiliser-isotypic reduction}\label{sec:reduction}

The multiplicities above are kernels of raising operators on the weight space
of $\lambda$, whose dimension $N_S$ grows quickly.  The following reduction,
which generalises the rectangular case of \cite{Companion}, is what makes the
computations feasible.

\begin{theorem}[Reduction]\label{thm:reduction}
Let $\lambda$ have equal-part blocks $B$ of sizes $k_B$ and values $m_B$, so
that the stabiliser $\Stab_W(\lambda)=\prod_B S_{k_B}$ permutes variable indices
within blocks and acts on the weight space of $\lambda$.  Every highest-weight
vector $v$ of weight $\lambda$ is an eigenvector,
$P_w\,v=\chi_\lambda(w)\,v$ with $\chi_\lambda(w)=\prod_B\operatorname{sgn}(w_B)^{m_B}$.
Hence the highest-weight vectors lie in the single $\chi_\lambda$-isotypic
component of the weight space, whose dimension we denote $n_\chi$.
\end{theorem}

\begin{proof}[Proof sketch]
$P_w v$ has weight $w(\lambda)=\lambda$ and lies in the irreducible generated
by $v$, whose $\lambda$-weight space is one-dimensional; the scalar is
intrinsic to the module, and on the leading-minor model a transposition inside
a block of value $m$ acts by $(-1)^{m}$.  The rectangular case is the special
value $\chi=\operatorname{sgn}^{\delta}$.  Note that the one-operator shortcut
available in the rectangular case relies on $2$-transitivity of the alternating
group and does not generalise; only the isotypic restriction does.
\end{proof}

The reduction was validated by computing, at seven cells, the full kernel
against every candidate character of the stabiliser and confirming it lands
entirely in $\chi_\lambda$ and is empty elsewhere --- a test that fails for a
wrong sign rule, and passes in every direction.

The dimension $n_\chi$ is an upper bound for the number of highest-weight
vectors, not a multiplicity or a rank, and it must be computed: it is given
exactly by the signed orbit count
$n_\chi=|\Stab_W(\lambda)|^{-1}\sum_{w}\chi_\lambda(w)\,\mathrm{Fix}(w)$, with
$\mathrm{Fix}(w)$ the number of weight-$\lambda$ monomials fixed by $w$.  The
quotient $N_S/|\Stab_W(\lambda)|$ is neither an upper nor a lower bound for
$n_\chi$: measured cells lie on both sides of it.  As an example of the
reduction, the unique degree-six invariant $I_6$ of senary quartics, of weight
$(4^{6})$, is reduced from $N_S=1\,532\,489$ to $n_\chi=2804$.  (The feasible
range is a property of the implementation rather than of the theorem; one
six-row cell since closed in this programme, $((12,4,4,4,4,4),8)$, has
$N_S=27\,009\,659$.)

\section{Negative results}\label{sec:negative}

We collect what the computations rule out.  All are det-side or reducible-side
statements.  By the washout theorem a length-five certificate would not see the
permanent's own structure: at $r=5$, $\Pp_5=\Rr_5$, so a positive gap there
would separate reducibility, not the permanent, from the determinant.  It would
nevertheless be a genuine multiplicity obstruction for the padded permanent,
precisely because $\Pp_5=\Rr_5$; the washout changes the interpretation of
such a certificate, not its existence, and is not an impossibility theorem.
The length-six data is the first that can see the permanent itself.

For a weight $\lambda$ with $\ell(\lambda)=k$ and a degree $\delta$ write
\[
  i_{\det}(\lambda,\delta)=a(\lambda,\delta)-\mult_\lambda\bC[\Ddet_k]_\delta ,
\]
the number of copies of $S_\lambda(\bC^k)$ in $I(\Ddet_k)_\delta$ (the
\emph{ideal-copy count}), as distinct from the coordinate-ring multiplicity
$\mult_\lambda\bC[\Ddet_k]_\delta$; and let
$\Delta(\lambda,\delta)=\mult_\lambda\bC[\Pp_k]_\delta-\mult_\lambda\bC[\Ddet_k]_\delta$
be the gap of Theorem~\ref{thm:transfer}.  By \eqref{eq:lengthred} both
multiplicities are those of the orbit closures in $\Sym^{4}\bC^{16}$.

\begin{theorem}[Blindness on the short slab]\label{thm:slab}
\begin{enumerate}
\item For every weight with $\ell(\lambda)\le4$ and every degree $\delta$,
$\Delta(\lambda,\delta)\le0$.
\item For every weight with $\ell(\lambda)\le3$ and every degree $\delta$,
$i_{\det}(\lambda,\delta)=0$, i.e.\
$\mult_\lambda\bC[\Ddet_{\ell(\lambda)}]_\delta=a(\lambda,\delta)$.
\item Let $e$ be the degree of a generator of the principal ideal
$I(\Ddet_4)\subset\bC[\Sym^{4}\bC^{4}]$.  For every weight with
$\ell(\lambda)=4$, $i_{\det}(\lambda,\delta)=0$ for $\delta\le e-1$.  Here
$e\ge10$ is certified (below), and $e=320112$ is adopted from
\cite[Thm.~2]{LLV}.  Hence $i_{\det}=0$ on the whole slab
$\{\ell(\lambda)\le4,\ \delta\le9\}$ unconditionally, and on
$\{\ell(\lambda)\le4,\ \delta\le e-1\}$ given the adopted value of $e$.  At
$\delta=e$ the generator itself is a copy of $S_{(e^{4})}$ in the ideal, so
$i_{\det}((e^{4}),e)\ge1$ and the equality $\mult=a$ of item~(2) does not
extend to length four in every degree.
\item On the range of items (2)--(3),
$\Delta(\lambda,\delta)=-\big(a(\lambda,\delta)-\mult_\lambda\bC[\Rr_k]_\delta\big)$,
$k=\ell(\lambda)$: a quantity about the reducible locus only.  It is strictly
negative at $((8,8,8),6)$ and at $((12,8,8),7)$, where $\Delta=-1$.
\end{enumerate}
\end{theorem}

\begin{proof}
(1) By \eqref{eq:lengthred} a weight of length $k\le4$ sees the two varieties
$\Pp_k$ and $\Ddet_k$.  For $k\le4$, $\Pp_k=\Rr_k$ (Theorem~\ref{thm:washout})
and $\Rr_k\subseteq\Ddet_k$ (Proposition~\ref{prop:contain}).  Containment of
varieties gives a surjection $\bC[\Ddet_k]\twoheadrightarrow\bC[\Pp_k]$ of
$\GL_k$-algebras, hence $\mult_\lambda\bC[\Pp_k]_\delta\le
\mult_\lambda\bC[\Ddet_k]_\delta$.  No degree enters: this item holds at every
$\delta$, including $\delta\ge e$.

(2) $\Ddet_3=\Sym^{4}\bC^{3}$ (the codimension table), and for $k\le2$
$\Ddet_k=\Sym^{4}\bC^{k}$ as well (every binary quartic is a product of linear
forms, the determinant of a diagonal matrix).  So $I(\Ddet_k)=0$ for $k\le3$
and $\mult=a$ at every weight of length at most three and every degree.

(3) $\Ddet_4$ is irreducible, as the closure of the image of an irreducible
variety, and has dimension $34$ in the $35$-dimensional $\Sym^{4}\bC^{4}$ (the
codimension table), so it is a hypersurface.  Its ideal in the factorial ring
$\bC[\Sym^{4}\bC^{4}]$ is principal, generated by an irreducible $g$ spanning a
$\GL_4$-stable line: a semi-invariant of weight $(e^{4})$, $e=\deg g$.  Hence
$I(\Ddet_4)_\delta=0$ for $\delta<e$, in every weight, and $g$ itself is the
copy asserted at $\delta=e$.  For $e\ge10$: at the rectangular weights
$((\delta^{4}),\delta)$ the multiplicity equals $a$ for $\delta=4,6,7,8$
(certified at both primes: a modular rank attaining $a$ is a certificate over
$\bQ$), while $a=0$ at the remaining $\delta\le9$; a generator of degree
$\le9$ would be a rectangular semi-invariant of that degree vanishing on
$\Ddet_4$, contradicting one of these.  The value $e=320112$ is the degree of
the component $F_1$ of \cite[Thm.~2]{LLV}, the closure of the quartic surfaces
with a $4\times4$ linear determinantal representation (admissible pair
$(5,5,5,5),(6,6,6,6)$ in \cite[Table~2]{LLV}); it is adopted, not certified
here, and nothing in this paper reaches degree $320112$.

(4) On that range $\mult_\lambda\bC[\Ddet_k]=a$, and $\Pp_k=\Rr_k$ for
$k\le4$.  At the two named cells, both of length three, the reducible locus
$\Rr_3$ carries exactly one ideal copy (proved and measured at both primes; at
$((8,8,8),6)$, where $a=2$, it is the combination of the square of the
degree-three invariant of ternary quartics and the degree-six invariant that
vanishes on every quartic with a linear factor), while
$i_{\det}=0$ by item~(2).
\end{proof}

In every degree, then, what holds on the length-$\le4$ slab is the inequality
of item~(1), a containment statement; the equality $\mult_\lambda\bC[\Ddet]=a$
holds at length at most three in every degree and at length four only below
the onset $e$.

\begin{theorem}[No arithmetic obstruction]\label{thm:arith}
Call a cell $(\lambda,\delta)$ \emph{occurrence-forced} if the ambient
multiplicity exceeds the Peter--Weyl bound of the determinant orbit,
$a(\lambda,\delta)>m_{\det}(\lambda)$.  No length-five cell at $n=4$ is
occurrence-forced through degree ten (an exhaustive check of $2585$ cells), and
none of length six through ten in the computed range.  Consequently the onset
of $I(\Ddet_5)$ is a genuine multiplicity phenomenon
$\mult<a\le m_{\det}$, invisible to the arithmetic that pinned the $n=3$ ideal,
and no occurrence obstruction is arithmetically available at $n=4$.
\end{theorem}

The contrast with $n=3$ is regime, not accident: at $n=3$ the ideal was first
pinned at lengths eight and nine by the degenerate mechanism $m_{\det}=0$, a
different length regime; at length five the $n=3$ route was silent too.  We
emphasise that Theorem~\ref{thm:arith} does not close the obstruction question:
a multiplicity obstruction with $a\le m_{\det}$ on both sides, detectable only
by an actual rank drop, is invisible to the arithmetic screen and lives in the
length-six-and-up regime that the reduction of \S\ref{sec:reduction} has only
begun to open.  In every length-six cell measured through degree seven the
determinant's ideal is empty and the padded multiplicity equals the
reducible-locus multiplicity, so no permanent-specific equation has been found.

\begin{remark}[The gate]\label{rem:gate}
Earlier sweeps in this programme used the gate $a\ge2$, on the ground that
occurrence obstructions ($a=1$, $m_{\det}=0$) are closed by \cite{BIP}.  That
theorem requires $n\ge m^{25}$ and is silent at $(m,n)=(3,4)$; the correct gate
at small parameters is $a\ge1$, and the $a=1$ cells --- one of which is the
degree-five reducible invariant of \S\ref{sec:reducible} --- must be screened.
After Proposition~\ref{prop:contain} the length gate is $\ell(\lambda)\ge5$,
and after Corollary~\ref{cor:kl} the first-row gate is $\lambda_1\ge\delta$.
This is a cautionary instance of transporting an asymptotic theorem into the
small regime.
\end{remark}

The honest frame, following \cite{DIP}: the set-theoretic non-containment
$x_0\per_3\notin\overline{\GL_{16}\cdot\det_4}$ is known (from
\cite[Thm.~1.0.1]{LMR}: $\overline{dc}(\per_m)\ge m^{2}/2$, so
$\overline{dc}(\per_3)\ge5>4$), so we are exhibiting the multiplicity method on
a separation established by other means.  What the present study contributes
is a map of where the certificate can live, with each coordinate labelled.
\emph{Length $\ge5$}: proved, since every weight of length at most four has
$\Delta\le0$ in every degree (Theorem~\ref{thm:slab}(1)); length five is
\emph{not} excluded (see the opening of this section).  \emph{First row
$\ge$ degree}: proved (Corollary~\ref{cor:kl}).  \emph{Degree $\ge8$}:
measured, not proved --- a positive gap needs $i_{\det}\ge1$, and the
determinant's five-row ideal $I(\Ddet_5)$ is measured empty through degree
seven, as is the determinant's ideal on every length-six cell measured through
degree seven; the coordinate holds only as far as those measurements reach.
Inside this map the open region is large: the $n=4$ region $\delta\le8$,
$5\le\ell(\lambda)\le\delta$, $\lambda_1\ge\delta$ contains exactly $4{,}198$
weight--degree labels (a certified count), $2{,}734$ of them with $a>0$, and
$2{,}571$ of those were still open at the time of writing.

\subsection*{The method exhibited on a separation of its own size.}

That frame invites an obvious test, and it can now be reported.  Take the
$n=3$ member of the family that supplies the length-nine equation, namely
$\lambda=(19,7,2^{5})$, $\delta=12$, $r=7$, with ambient multiplicity $a=6$; and
compare $\det_3$ against the \emph{unpadded} $\per_3$, both cubics in nine
variables restricted to seven.  The comparison is favourably oriented, unlike
the padded one at small length: writing $\Ddet_7$ and $\Pp_7$ for the two
varieties of cubics in $\Sym^{3}\bC^{7}$ (of dimension $84$),
\[
  \dim\Ddet_7 = 9r-16 = 47 ,\qquad \dim\Pp_7 = 9r-4 = 59 ,
\]
so the determinant variety is the smaller of the two and neither fills the
ambient space.  We find
\[
  \mult\nolimits_{\det}=5,\qquad \mult\nolimits_{\per}=6=a ,\qquad
  \Delta = \mult\nolimits_{\per}-\mult\nolimits_{\det} = +1 ,
\]
proved; its one premise beyond the measurements, $(\star)$ below, is itself proved.  The halves are rigorous for
different reasons, and it is worth separating them.  $\mult_{\per}=a$ was
measured as a rank modulo $p$ at both house primes, and
$\rank_p\le\rank_\bQ\le a$ forces equality over $\bQ$; there is no
finite-point caveat in this direction, since the caveat is that vanishing at
finitely many points does not prove ideal membership and here nothing vanishes.
The same argument applied to the determinant's measured rank $5$ gives
$\mult_{\det}\ge5$.  The reverse inequality is an ideal-membership statement,
and it comes from \cite{LMR}: \cite[Thm.~2.3.1]{LMR} with \cite[Thm.~3.1.1 and
\S3.2]{LMR} put a copy of the module of highest weight
$12\omega_1+5\omega_2+2\omega_7=(19,7,2^{5})$ in the ideal of
$\overline{\GL_9\cdot\det_3}$ in degree $12$ (the example printed in
\cite[\S3.2]{LMR}, which also records the ambient multiplicity six).  So
$i_{\det}\ge1$ in nine variables, and $\mult_{\det}\le5$ once that copy is
transferred to the seven-variable model $\Ddet_7$.  The transfer is the
programme's isotypic restriction from nine variables to seven, valid for
$\ell(\lambda)\le7$ and used here at its boundary $\ell(\lambda)=7$; call it
$(\star)$.  It is not a statement of \cite{LMR}; it is proved in this programme, so the conclusion
$\mult_{\det}=5$, hence $\Delta=+1$, is proved.  (We use only
the half of the sentence of \cite[\S3.2]{LMR} that is proved there,
$i_{\det}\ge1$; the bound $i_{\det}\le1$ comes from the measured rank.)  Both
multiplicities are nonzero, so
$S_\lambda$ occurs in both coordinate rings: this is a \emph{multiplicity}
obstruction and not an occurrence obstruction --- which is the stronger reading,
since \cite{BIP} closes the occurrence route in the padded regime and leaves the
multiplicity route open.

We claim nothing new for the separation itself: $\dim\Pp_7=59>47=\dim\Ddet_7$
already gives $\Pp_7\not\subseteq\Ddet_7$ by a dimension count.  The point is
the certificate.  It is the first $\Delta>0$ this programme has produced end to
end by its own machinery, at a cell of $17\,047$ coordinates rather than the
$3.1\times10^{7}$ of the length-nine cell, and it exercises exactly the pipeline
the length-nine comparison requires: one source, one raising kernel, two
evaluation families, two primes.  It is unpadded and at $n=3$, so it bears on the
padded model only as calibration --- indeed the \emph{padded} $n=3$ cell can
never carry an obstruction, since $\per_2$ and $\det_2$ are both rank-four
quadratic forms in four variables and hence $\GL$-equivalent, so the padded
$\per_2$ in degree three is $\ell\cdot\det_2$ up to $\GL$, which is the block
split $\operatorname{diag}(\ell,A_2)$ inside $\det_3$ and forces
$\Delta\le0$ there identically.

\subsection*{What the length-nine cell reduces to.}

At $n=4$ the smallest length carrying a non-vacuous determinant equation is
$r=9$, and the deduction there needs neither the target dimension nor
stability.  The ambient multiplicities along the ladder
$(4\delta-31,17,2^{7})$ end $\dots,272,273,274$, so the last two increments are
one apiece, and the ladder's monotonicity turns the whole question into a single
rank lower bound: $\rank\Theta^{+}\ge273$ at the predecessor forces
$\mult_{\det}\ge273$ at the goal cell, hence $i_{\det}\le1$, hence
$i_{\det}=1$ given $i_{\det}\ge1$.  The last is \cite{LMR}: by
\cite[Thm.~2.3.1]{LMR} the variety $\mathrm{Dual}_{k,d,N}$ has equations
forming a copy of the module of highest weight
$\Omega(k,d)=(d-1)(d-2)(k+2)\omega_1+(d(k+2)-2k-5)\omega_2+2\omega_{k+3}$, in
degree $(k+2)(d-1)$, and by \cite[Thm.~3.1.1]{LMR} $[\det_n]$ lies on
$\mathrm{Dual}_{2n-2,n,n^{2}}$; at $n=d=4$, $k=6$, $N=16$ this is the weight
$48\omega_1+15\omega_2+2\omega_9=(65,17,2^{7})$ in degree $24$, and the
length reduction \eqref{eq:lengthred} carries the copy from sixteen variables to
$\Ddet_9$.  That transfer is the $n=4$ analogue of the premise $(\star)$ of the
$n=3$ control above: the $n=3$ case is proved; the $n=4$ case is not yet checked, so
this transfer is flagged, not established.  (We do not
use \cite[Thm.~1.0.2]{LMR}, whose printed coefficient of $\omega_1$ and degree
are both half of these values.)  Equivalently, in the
room-one form, the equation is born exactly when the Schur complement
$c-b^{\mathsf T}A^{-1}b$ vanishes, with $A$ the degree-$24$ Gram restricted to
the transported predecessor.  The $n=3$ shadow above exhibits the same mechanism
at a size where it can be run.

That rank lower bound has since been established.  A basis of
$M_{\lambda}=\mathrm{HWV}_{\lambda}(\Sym^{24}\Sym^{4}\bC^{9})$ of all $274$
vectors was constructed by births along the ladder, each rung's birth quotient
$M_d/uM_{d-1}$ certified by a nonzero minor at both primes, where
$u=c_{(4,0,\dots,0)}$; multiplication by $u$ is injective and
$\ker(M_d\to R/(u))=uM_{d-1}$, so the transported births are a basis, which the
generic column confirms.  Evaluating at determinant pencils gives a nonzero
$273\times273$ minor \emph{on the transported degree-$23$ rows alone}, so
$\Theta^{+}$ is injective on $uM_{23}$ and $i_{\det}(23)=0$; with
$b_{24}=1$ and \cite{LMR} this pins
\[
  \rank\Theta_{\lambda}^{\det}=273,\qquad i_{\det}=1
\]
exactly, and no degree-$24$ candidate enters the argument.  The one-dimensional
kernel is therefore \cite{LMR}'s own equation, and its residues are nonzero at
every one of $282$ integer padded points at both primes, so it does not vanish on
the padded permanent.

On the padded side a nonzero $269\times269$ minor gives the floor
$\rank\Theta_{\lambda}^{\mathrm{pad}}\ge269$, hence $i_{\mathrm{pad}}\le5$ and
$\Delta=1-i_{\mathrm{pad}}\ge-4$.  The observed nullity of five is a ceiling read
off a finite point set and is not a theorem: a nonzero minor bounds $i$ from
above, and only a membership proof bounds it from below.  Such a proof is
available.  At the ladder rung $(21,17,2^{7})$ in degree $13$ the reducible
locus carries exactly three ideal copies: evaluation at the chosen points is
injective on the relevant $73$-dimensional source space, certified by a
nonzero $73\times73$ minor at both primes, so the three sampled kernel
directions are genuine equations of $\Rr_9$.  Multiplication by the
eleventh power of $u$ carries them injectively to three ideal copies at the
goal cell, and $\Pp_9\subseteq\Rr_9$ gives $i_{\mathrm{pad}}\ge3$.  Hence
\[
  \Delta\in[-4,-2],
\]
certified conditional on the adopted value $73$ of the dimension of that source
space (adopted on two independent lineages, not proved here); in particular
$\Delta=+1,0,-1$ are excluded at this cell.  Two further facts are measured on
the sampled kernel, at both primes.  None of the five sampled
kernel directions meets the degree-$24$ birth line, so
$i_{\mathrm{pad}}(24)=i_{\mathrm{pad}}(23)$ and $\Delta=1-i_{\mathrm{pad}}(23)$;
and the padded and reducible kernels coincide, so
$\mult_{\mathrm{pad}}=\mult_{\mathrm{red}}$ and the whole question is one about
$\Rr_9$ rather than about the permanent.  We record, with the caveat that any
exhibited element of $I(\Ddet)^{\mathrm{HWV}}$ obtained by evaluation carries:
vanishing at finitely many points is Schwartz--Zippel evidence of ideal
membership, not a proof of it, and the rigorous lower bound $i_{\det}\ge1$ on
the determinant side remains \cite{LMR}'s (Theorems~2.3.1 and~3.1.1 there).

\section{Open problems}\label{sec:open}

\begin{question}
Is Conjecture~\ref{conj:onset} true at $n=3$?  The decisive computations are the
balanced length-five cells at degree eight and nine, now within reach of the
reduction of \S\ref{sec:reduction}; each empty cell is evidence, and a single
bite below degree $65$ settles it in the negative.
\end{question}

\begin{question}
Does a multiplicity obstruction between $\Pp_6$ and $\Ddet_6$ exist at length
six?  By Proposition~\ref{prop:pieri} the first candidate degree is seven.  The
$682$ six-row cells we have measured through degree twelve --- the skewed ones,
which are the affordable ones --- carry no determinant equation, no
permanent-specific equation and no $\per_4$ equation (measured, not proved),
and $63$ tails are closed
at every degree by the first-stable-cell argument.  The balanced six-row cells,
with $n_\chi\ge10^{6}$, remain unmeasured and are the frontier.
\end{question}

\begin{question}
Is there a closed form for the extra Jacobian syzygy of a determinantal
threefold --- the one relation beyond Koszul that Step 3 of
Theorem~\ref{thm:cap} produces?  An explicit adjugate identity would make the
cap theorem elementary and remove the appeal to Dimca and Gulliksen--Neg\aa rd.
\end{question}

\begin{question}
Does the reducible-locus screen (Theorem~\ref{thm:transfer}) ever certify
blindness over an unbounded degree range?  Direction 1's Kempf-collapsing
computation of $\bC[\Rr_r]$ weight by weight, if carried out, would decide
whether $\Delta\le0$ holds for all $\delta$ on a fixed length.
\end{question}

\begin{question}
Is $I(D^{\per_3}_r)_\delta$ ever nonzero?  By
Proposition~\ref{prop:pieri} this single cubic-side ideal governs the whole
question of whether the permanent is visible at all: it is empty in degree
$\delta$ at length $r$ if and only if $\mult_\lambda\bC[\Pp_r]=
\mult_\lambda\bC[\Rr_r]$ at \emph{every} weight of that length and degree.  It
is zero for every $r$ in every degree $\delta\le8$, and at $r=6$ in degree
nine (\S\ref{sec:intro}); its constituents have length between $6$ and
$\min(r,\delta)$, with no further upper bound on that length known, and the
first open cases lie in degree nine.  One scan settles a whole
degree, so
this is the cheapest open question the programme has, and a nonzero weight would
be the first place the permanent becomes visible to a multiplicity.
\end{question}

\section*{Acknowledgements}

The work reported here was carried out using AI tools, including Claude (Anthropic).
Responsibility for the correctness of every statement in this paper rests with
the author.

\section*{Data availability}\label{sec:data}

All code, inputs, exact outputs and the full commit history --- including the
pre-registrations of the predictions tested --- are available at
\url{https://github.com/swsethuraman/gct}.  The measured quantities of
\S\ref{sec:cap} and \S\ref{sec:negative} are reproduced by the scripts there
over two word-size primes; the symbolic identities of \S\ref{sec:cap} are
checked in exact rational arithmetic.

\begin{thebibliography}{99}

\bibitem{Beau} A.~Beauville, \emph{Determinantal hypersurfaces}, Michigan Math.
J. \textbf{48} (2000), 39--64.

\bibitem{BrunsVetter} W.~Bruns, U.~Vetter, \emph{Determinantal Rings}, Lecture
Notes in Math. \textbf{1327}, Springer, 1988.

\bibitem{BI} P.~B\"urgisser, C.~Ikenmeyer, \emph{Fundamental invariants of
orbit closures}, J. Algebra \textbf{477} (2017), 390--434.

\bibitem{BIP} P.~B\"urgisser, C.~Ikenmeyer, G.~Panova, \emph{No occurrence
obstructions in geometric complexity theory}, J. Amer. Math. Soc. \textbf{32}
(2019), 163--193.

\bibitem{Companion} S.~Sethuraman, \emph{Conductors of orbit closures, and the
fundamental invariant of the $3\times 3$ determinant}, arXiv preprint, 2026,
arXiv:XXXX.XXXXX [identifier to be inserted at submission].

\bibitem{Companion2} S.~Sethuraman, \emph{Singular matrix spaces and the
length-five non-containment}, technical report, this programme, 2026; PDF at
[repository URL to be inserted at submission].

\bibitem{Companion3} S.~Sethuraman, \emph{$D_{45}\cap P_5$ classified up to its
boundary, and the rank thresholds at $N=6,7,8$}, technical report, this programme,
2026; PDF at [repository URL to be inserted at submission].

\bibitem{Dimca13} A.~Dimca, \emph{Syzygies of Jacobian ideals and defects of
linear systems}, Bull. Math. Soc. Sci. Math. Roumanie \textbf{56(104)} (2013),
191--203; arXiv:1210.1795.

\bibitem{DIP} J.~D\"orfler, C.~Ikenmeyer, G.~Panova, \emph{On geometric
complexity theory: multiplicity obstructions are stronger than occurrence
obstructions}, ICALP 2019 / SIAM J. Appl. Algebra Geom.

\bibitem{GN} T.~H.~Gulliksen, O.~G.~Neg\aa rd, \emph{Un complexe r\'esolvant
pour certains id\'eaux d\'eterminantiels}, C. R. Acad. Sci. Paris \textbf{274}
(1972), 16--18.

\bibitem{KL2} E.~Kadish, J.~M.~Landsberg, \emph{Padded polynomials, their
cousins, and geometric complexity theory}, Comm. Algebra \textbf{42} (2014),
2171--2180; arXiv:1204.4693.

\bibitem{Kleiman} S.~L.~Kleiman, \emph{The transversality of a general
translate}, Compositio Math. \textbf{28} (1974), 287--297.

\bibitem{Landsberg} J.~M.~Landsberg, \emph{Geometry and Complexity Theory},
Cambridge Stud. Adv. Math. \textbf{169}, Cambridge Univ. Press, 2017.

\bibitem{LLV} M.~Leal, C.~Lozano Huerta, M.~Vite, \emph{The Noether--Lefschetz
locus of surfaces in $\mathbb{P}^3$ formed by determinantal surfaces}, Math.
Nachr. \textbf{297} (2024), 4671--4688; arXiv:2303.09028.

\bibitem{LMR} J.~M.~Landsberg, L.~Manivel, N.~Ressayre, \emph{Hypersurfaces
with degenerate duals and the geometric complexity theory program}, Comment.
Math. Helv. \textbf{88} (2013), 469--484.

\bibitem{MS} K.~D.~Mulmuley, M.~Sohoni, \emph{Geometric complexity theory I},
SIAM J. Comput. \textbf{31} (2001), 496--526.

\bibitem{Valiant} L.~G.~Valiant, \emph{Completeness classes in algebra}, Proc.
11th ACM STOC (1979), 249--261.

\end{thebibliography}

\end{document}
